\SourceLeaf{372}%
Es ist kaum zu begreifen, wie er sich auf diese Weise selbst täuschen konnte. Er erhebt sogar selbst den Einwand, dass die Summe der Stücke \(Sp,pp'\) und so weiter, wenn die Stücke immer kleiner werden, trotz ihrer unaufhörlich wachsenden Anzahl über eine gewisse Grenze nicht hinauswachsen könnte. Er meint, den Einwand dadurch entkräften zu können, dass jene Stücke, auch wenn sie immer kleiner werden, doch stets größer bleiben als \textit{eine angebbare Größe}. Denn diese Stücke sind Katheten rechtwinkliger Dreiecke und folglich immer größer als der Unterschied zwischen der Hypotenuse und der anderen Kathete. Fast scheint es, als hätte eine grammatische Zweideutigkeit den Verfasser in die Irre geführt, nämlich der doppelte Sinn des Artikels \textit{eine} in eine angebbare Größe. Der Schluss des Verfassers wäre nur dann richtig, wenn sich zeigen ließe, dass die Stücke \(Sp,pp'\) und so weiter stets größer bleiben. als \textit{eine bestimmte} angebbare Größe, zum Beispiel als der Unterschied zwischen der Hypotenuse \(pT\) und der Kathete \(ST\). Das lässt sich aber nicht beweisen. Es lässt sich nur zeigen, dass jedes Stück stets größer bleibt als eine angebbare Größe, die für jedes Stück eine andere ist. So ist \(Sp\) größer als der Unterschied zwischen \(pT\) und \(ST\), ferner \(pp'\) größer als der Unterschied zwischen \(qp'\) und \(qp\) und so weiter. Damit verliert der Beweis seine ganze Kraft.

Auf dieselbe Weise, auf die er seinen Beweis führen zu können glaubte, könnte er auch beweisen, dass in einem ebenen Dreieck \(ABC\), in dem \(B\) ein rechter Winkel ist, \(C\) nicht spitz sein könne. Er müsste nur von \(B\) die Senkrechte \(BD\) auf die Hypotenuse \(AC\) fällen, dann die Senkrechte \(DE\) auf \(AB\) und so ohne Ende die Senkrechten \(EF,FG,GH\) und so weiter abwechselnd auf \(AC\) und \(AB\). Die Stücke \(CD,DF,FH\) und so weiter sind stets größer als der angebbare Unterschied zwischenHypotenuse und einer Kathete \textit{desjenigen} rechtwinkligen Dreiecks, in dem jeweils das betreffende Stück die andere Kathete ist. Trotzdem schöpft ihre Summe die Hypotenuse \(AC\) offenbar nie aus, wie groß man ihre Anzahl auch wählt.

Wir müssten fast bedauern, bei so bekannten und einfachen Dingen so lange verweilt zu haben, wenn diese Schrift nicht durch die Art, wie sie schon vor ihrem Erscheinen in öffentlichen Blättern angekündigt wurde, mehr Aufmerksamkeit als gewöhnlich auf sich gezogen hätte. Dem Verfasser scheint es übrigens wirklich um die Wahrheit zu gehen. Wir bemerken daher nur noch, dass der Verfasser anschließend auf ganz ähnliche und damit ebenso untaugliche Weise beweisen %
